\chapter{Longueurs, angles et lieux géométriques}
\section*{Exercice 17.1 -- Al-Kashi}
Un triangle a deux côtés de longueurs $4$ et $6$, formant un angle $120^\circ$. Calculer le troisième côté.

\section*{Exercice 17.2 -- Milieu et lieu}
$AB=6$ et $I$ est son milieu. Décrire les lieux $\overrightarrow{MA}\cdot\overrightarrow{MB}=k$ pour $k=0,-9,-10$.

\section*{Exercice 17.3 -- Médiane}
Dans un triangle, $AB=5$, $AC=7$, $BC=8$. Calculer la longueur de la médiane issue de $A$.

\section*{Exercice 17.4 -- Aire par un angle}
Deux côtés de longueurs $5$ et $8$ forment un angle $\pi/6$. Calculer l'aire en justifiant la formule.

\section*{Exercice 17.5 -- Centre de gravité}
Pour $A=(0,0)$, $B=(6,0)$, $C=(0,3)$, déterminer $G$ et prouver qu'il minimise $MA^2+MB^2+MC^2$.

\section*{Exercice 17.6 -- Hauteurs sans dessin}
Pour $A=(0,0)$, $B=(4,0)$, $C=(1,3)$, calculer l'orthocentre et vérifier la troisième hauteur.
